\documentclass[11pt]{article}
\usepackage[T1]{fontenc}
\usepackage{lmodern}
\usepackage[margin=1in]{geometry}
\usepackage{amsmath,amssymb,amsthm}
\usepackage{booktabs}
\usepackage{tabularx}
\usepackage{microtype}
\usepackage[title]{appendix}
\usepackage[colorlinks=true,linkcolor=blue,citecolor=blue,urlcolor=blue]{hyperref}
\hypersetup{pdftitle={A Steiner triple system with a broken chromatic spectrum},pdfauthor={P. Graham}}
\newtheorem{theorem}{Theorem}
\newtheorem{corollary}[theorem]{Corollary}
\newtheorem{lemma}[theorem]{Lemma}
\theoremstyle{definition}
\newtheorem{definition}{Definition}
\title{A Steiner triple system with a broken chromatic spectrum}
\author{P. Graham\\ \small\url{https://orcid.org/0009-0007-8493-1505}}
\date{October 31, 2026}

\begin{document}
\maketitle
\begin{abstract}
A bicoloring of a Steiner triple system colors its points so that every
block receives exactly two colors, and the feasible set of the system is
the set of integers $k$ for which some bicoloring uses exactly $k$ colors.
Mario Gionfriddo conjectured in 2003 that the feasible set of every
bicolorable Steiner triple system has no gaps; Tuza and Voloshin listed the
question as open in 2008, and Gionfriddo restated the conjecture in 2025.
In 2014 Gionfriddo, Guardo, Milazzo and Voloshin left open whether a
system of order 39 has feasible set $\{3,5\}$. We give a Steiner triple
system of order 39 that has bicolorings with exactly three and exactly
five colors but none with exactly four, so the conjecture is false. With
their theorem that bicolorings of systems of order 39 use three, four or
five colors, its feasible set is $\{3,5\}$. The class sizes of the two
bicolorings come from the tables of Colbourn, Dinitz and Rosa. The
absence of a four-color bicoloring is proved by computer: a SAT solver's
clausal refutation of a direct encoding, checked by DRAT-trim and by a
separate unit-propagation checker.

This work has not yet been reviewed by experts.
The proofs, the search code, the certificate checkers and the text were
produced by AI models (from OpenAI and Anthropic), working under the
author's direction.
\end{abstract}
\medskip

\noindent\textit{Mathematics Subject Classification:} 05B07, 05C15, 05C65.
\textit{Keywords:} Steiner triple system, bicoloring, mixed hypergraph,
feasible set, chromatic spectrum.

\section{Introduction}

A \emph{Steiner triple system} of order $v$, an STS$(v)$, is a pair
$(X,\mathcal{B})$ where $X$ is a set of $v$ points and $\mathcal{B}$ is a
family of $3$-element subsets of $X$, the \emph{blocks}, such that every
pair of distinct points lies in exactly one block. Such a system exists
exactly when $v\equiv 1$ or $3 \pmod 6$. The smallest nontrivial example is
the Fano plane, with $7$ points and $7$ blocks.

\begin{definition}
Let $S=(X,\mathcal{B})$ be a Steiner triple system and $k$ a positive
integer. A \emph{$k$-bicoloring} of $S$ is a surjective map
$f\colon X\to\{1,\dots,k\}$ such that every block $B\in\mathcal{B}$
satisfies $|f(B)|=2$. The \emph{feasible set} $\Omega(S)$ is the set of
all $k$ for which $S$ has a $k$-bicoloring. A \emph{gap} of $\Omega(S)$ is
an integer $k\notin\Omega(S)$ with elements of $\Omega(S)$ both below and
above $k$.
\end{definition}

For example, coloring a point of the Fano plane $1$, the two other points
of one block through it $2$, and the remaining four points $3$ gives a
$3$-bicoloring, and the Fano plane has feasible set $\{3\}$.
Bicolorings are the strict colorings of the system viewed as a mixed
hypergraph in Voloshin's sense~\cite{Voloshin95}, in which every block is
both a C-edge (it has two points of the same color) and a D-edge (it has
two points of different colors). Since a block is never monochromatic,
each color class of a bicoloring is an independent set, and a Steiner
triple system with more than three points has no $1$- or
$2$-bicoloring~\cite[Section~1]{CDR}. The
feasible set of a mixed hypergraph can have gaps: Jiang, Mubayi, Tuza,
Voloshin and West showed this in \emph{The chromatic spectrum of mixed
hypergraphs}~\cite{JMTVW}.

For Steiner triple systems, Gionfriddo conjectured the opposite in a
lecture at Messina in 2003, published in~\cite{Gionfriddo03}. Tuza and
Voloshin's survey of open problems on mixed hypergraphs, in its Section~7
on Steiner systems, records that ``M. Gionfriddo conjectured in [21] that
every BSTS has a gap-free chromatic spectrum'' (their [21]
is~\cite{Gionfriddo03}), where a BSTS is a Steiner
triple system whose blocks are both C-edges and D-edges, and lists ``Prove
or disprove that all BSTS have gap-free chromatic spectrum'' as Open
problem~2~\cite{TV08}. Gionfriddo's 2025 survey states the conjecture
again, as Conjecture [M. Gionfriddo 2003] in its
Section~4~\cite{Gionfriddo25}, in the terminology of chromatic spectra:
the spectrum is \emph{broken} when the feasible set has a gap. Gionfriddo, Guardo, Milazzo and
Voloshin~\cite{GGMV14} wrote that ``The problem to find a bicolorable STS
with a gap in its feasible set still remains an open problem'', listed
``Determine whether there exist systems $S$ with a gap in their
$\Omega(S)$'' as problem~4 of their conclusion, and, in the table closing
their Section~3, listed $\Omega(S)=\{3,5\}$ for an STS$(39)$ as open.

\subsection{Main results}

Let $S_{39}$ be the Steiner triple system on the points $0,1,\dots,38$
whose $247$ blocks are listed in Appendix~\ref{app:data}.

\begin{theorem}\label{thm:main}
The Steiner triple system $S_{39}$ has a $3$-bicoloring and a
$5$-bicoloring, and it has no $4$-bicoloring. Hence $4$ is a gap of
$\Omega(S_{39})$, and Gionfriddo's conjecture is false.
\end{theorem}

The bicolorings in Theorem~\ref{thm:main} have color classes of sizes
$10,12,17$ and $1,2,4,16,16$. The nonexistence of a $4$-bicoloring is
proved by computer, by a clausal proof of unsatisfiability checked by two
independent programs (Section~\ref{sec:verification}).

Gionfriddo, Guardo, Milazzo and Voloshin~\cite[Theorem~3.7]{GGMV14} state
that the union of the feasible sets of all STS$(39)$s is $\{3,4,5\}$.
Together with Theorem~\ref{thm:main} this gives the following.

\begin{corollary}\label{cor:omega}
$\Omega(S_{39})=\{3,5\}$.
\end{corollary}

This is the entry left open for order $39$ in the table
of~\cite{GGMV14}. The analogous question for the union of the feasible
sets over all systems of one order was settled earlier: Colbourn, Dinitz
and Rosa found that this union is $\{3,5\}$ for order~$31$ and remarked
that it ``is not necessarily an interval''~\cite[Section~4]{CDR}, while
Gionfriddo, Guardo, Milazzo and Voloshin proved that no single STS$(31)$
has feasible set $\{3,5\}$~\cite[Theorem~3.1]{GGMV14}. Theorem~\ref{thm:main}
concerns one system with both colorings at once.

The class sizes we used come from Colbourn, Dinitz and Rosa~\cite{CDR}.
Their list of $3$-splits in Section~2 contains $(17,12,10)$ for $v=39$,
and their Table~2 of $5$-splits contains $(16,16,4,2,1)$, each marked as
realized by some STS$(39)$. These realizations are separate systems. The
STS$(39)$ with a $5$-bicoloring described in the proof
of~\cite[Theorem~3.7]{GGMV14} ``does not admit two distinct disjoint
independent sets of cardinalities 9, 14, 16 and 10, 12, 17, i.e., it is
not 3-bicolorable.'' What is new here is a single system that carries a
$3$-bicoloring and a $5$-bicoloring together while admitting no
$4$-bicoloring; the difficult step is the last one, a nonexistence
statement over all colorings of $39$ points.

\section{The construction}\label{sec:construction}

We looked for the system and its two bicolorings at the same time. Fix
the $5$-coloring $g$ of $\{0,\dots,38\}$ with the classes
\[
\{0\},\quad \{1,2\},\quad \{3,4,5,6\},\quad \{7,\dots,22\},\quad \{23,\dots,38\}.
\]
A Boolean variable for each triple of points receiving exactly two colors
under $g$ records whether the triple is a block; the constraints require
every pair of points to lie in exactly one chosen triple. Further
variables describe a second coloring $h$ with three classes of sizes
$10$, $12$ and $17$, and every chosen triple must receive exactly two
colors under $h$. Every solution is therefore a Steiner triple system with
the $5$-bicoloring $g$ and a $3$-bicoloring $h$. The SAT solver
CaDiCaL~\cite{CaDiCaL} returned a solution, and $S_{39}$ is that solution,
with $h$ the coloring of Appendix~\ref{app:data}.

Nothing in this search excludes a $4$-bicoloring. That is proved separately,
for the system found, by the argument of the next section.

\section{Proof of Theorem~\ref{thm:main}}\label{sec:proof}

We encode $k$-bicolorings as satisfying assignments of a propositional
formula. For a Steiner triple system $S=(X,\mathcal{B})$ with
$X=\{0,\dots,v-1\}$ and an integer $k\ge 1$, let $\Phi_k(S)$ be the formula
in the variables $p_{x,c}$ ($x\in X$, $1\le c\le k$), read ``$x$ has color
$c$'', with the clauses
\begin{align*}
&p_{x,1}\vee\dots\vee p_{x,k} \quad\text{and}\quad \neg p_{x,c}\vee\neg p_{x,c'}
  &&\text{for } x\in X \text{ and } c<c',\\
&p_{0,c}\vee p_{1,c}\vee\dots\vee p_{v-1,c} &&\text{for } 1\le c\le k,\\
&\neg p_{x,a}\vee \neg p_{y,b}\vee\neg p_{z,c} &&\text{for }
  \{x,y,z\}\in\mathcal{B} \text{ and colors } a,b,c \text{ with } |\{a,b,c\}|\in\{1,3\}.
\end{align*}
In the last line each block is written once as $\{x,y,z\}$, and
$(a,b,c)$ runs over the assignments of colors to $x,y,z$ that give the
block one or three distinct colors; for $k=4$ there are $4+24=28$ of them.

\begin{lemma}\label{lem:encoding}
The satisfying assignments of $\Phi_k(S)$ correspond bijectively to the
$k$-bicolorings of $S$, by $f(x)=c$ exactly when $p_{x,c}$ is true.
\end{lemma}

\begin{proof}
The clauses of the first line hold exactly when each point has exactly
one color, so they correspond to maps $f\colon X\to\{1,\dots,k\}$. Given
such a map, the clauses of the second line hold exactly when $f$ is
surjective. For a block $\{x,y,z\}$, the clauses of the third line hold
exactly when $(f(x),f(y),f(z))$ is not an assignment with one or three
distinct colors, that is, when $|f(\{x,y,z\})|=2$.
\end{proof}

\begin{proof}[Proof of Theorem~\ref{thm:main}]
The first two assertions are finite checks. The $247$ blocks of
Appendix~\ref{app:data} cover every pair of points exactly once, and the
two colorings given there are surjective and give every block exactly two
colors; both facts are verified directly from the definitions
(Section~\ref{sec:verification}).

For the third assertion take $k=4$. The formula $\Phi_4(S_{39})$ has $156$
variables and $39\cdot 7+4+247\cdot 28=7{,}193$ clauses, and by
Lemma~\ref{lem:encoding}, $S_{39}$ has a $4$-bicoloring if and only if
$\Phi_4(S_{39})$ is satisfiable. CaDiCaL produced a proof of
unsatisfiability of $\Phi_4(S_{39})$ in the DRAT format~\cite{DRATtrim}: a
sequence of clause additions and deletions, ending with the empty clause,
in which every added clause follows from the formula and the clauses kept
so far by unit propagation. This proof is checked as described in
Section~\ref{sec:verification}. Hence $\Phi_4(S_{39})$ is unsatisfiable and
$S_{39}$ has no $4$-bicoloring.
\end{proof}

The encoding contains no symmetry breaking, no fixed colors and no
assumption on class sizes, so the refutation covers every map from the
$39$ points onto four colors.

\section{Computational verification}\label{sec:verification}

The system, its two bicolorings, the formula $\Phi_4(S_{39})$ in DIMACS
format and the DRAT proof are distributed with this paper, together with
the programs that check them; the result's README gives the single
command that runs every check, in a few seconds on a laptop. The checks
are as follows.
\begin{enumerate}
\item A direct check that the $247$ blocks are distinct $3$-subsets of
  $\{0,\dots,38\}$ covering every pair exactly once, and that both colorings
  are surjective and give every block exactly two colors.
\item An independent construction of $\Phi_4(S_{39})$ from the blocks,
  in which the block clauses are produced by listing all $4^3$ color
  assignments of each block and keeping those with one or three colors.
  It is compared, clause by clause as a multiset, with the distributed
  formula.
\item The DRAT proof is checked against the formula by
  DRAT-trim~\cite{DRATtrim}. Every lemma in its verified core is a
  unit-propagation consequence, so the proof needs no
  resolution-asymmetric-tautology step.
\item The same proof is checked by a short separate program that accepts an
  added clause only when unit propagation from its negation, over the
  formula and the clauses accepted before it, yields a conflict, and that
  ignores deletions. It accepts all $6{,}374$ nonempty added clauses and
  the final empty clause. Ignoring deletions keeps the check sound.
\end{enumerate}
As calibration, the same definitions give feasible set $\{3\}$ for the
Fano plane and for the affine plane of order $3$, by listing all set
partitions of their points, in agreement with~\cite[table at the end of
Section~3]{GGMV14}. For the Fano plane and $k=2,3,4$, the formula
$\Phi_k$ holds at exactly the $k$-bicolorings among all $k^7$ colorings.
Both proof checkers reject a proof consisting only of the empty clause,
the first half of the proof, and the proof applied to the satisfiable
formula $\Phi_5(S_{39})$; the design checker rejects altered blocks and
altered colorings.

Beyond the correctness of these programs and of the C compiler and Python
interpreter running them, the proof relies on nothing further. The
system was found by search, but the search is not part of the proof.
Corollary~\ref{cor:omega} relies in addition on~\cite[Theorem~3.7]{GGMV14}.

\section*{Acknowledgments}
The question is Mario Gionfriddo's, and the colorings were guided by the
class-size tables of Charles Colbourn, Jeffrey Dinitz and Alexander Rosa
and by the study of small feasible sets by Mario Gionfriddo, Elena Guardo,
Lorenzo Milazzo and Vitaly Voloshin. The SAT solver CaDiCaL is by Armin
Biere and collaborators; DRAT-trim is by Nathan Wetzler, Marijn Heule and
Warren Hunt.

\begin{thebibliography}{99}

\bibitem{CaDiCaL}
A.~Biere, T.~Faller, K.~Fazekas, M.~Fleury, N.~Froleyks and F.~Pollitt,
\emph{CaDiCaL 2.0}, in: Computer Aided Verification (CAV 2024), Lecture
Notes in Comput. Sci. 14681, Springer, 2024, 133--152.
\href{https://doi.org/10.1007/978-3-031-65627-9_7}{doi:10.1007/978-3-031-65627-9\_7}.

\bibitem{CDR}
C.~J. Colbourn, J.~H. Dinitz and A.~Rosa, \emph{Bicoloring Steiner triple
systems}, Electron. J. Combin. \textbf{6} (1999), \#R25.
\href{https://doi.org/10.37236/1457}{doi:10.37236/1457}.

\bibitem{Gionfriddo03}
M.~Gionfriddo, \emph{Colourings of hypergraphs and mixed hypergraphs},
Rend. Sem. Mat. Messina Ser.~II \textbf{9} (2003), 87--97.

\bibitem{Gionfriddo25}
M.~Gionfriddo, \emph{Some problems and conjectures about Voloshin triple
systems}, Mathematics \textbf{14} (2026), no.~1, 42, published
December 22, 2025.
\href{https://doi.org/10.3390/math14010042}{doi:10.3390/math14010042}.

\bibitem{GGMV14}
M.~Gionfriddo, E.~Guardo, L.~Milazzo and V.~Voloshin, \emph{Feasible sets
of small bicolorable STSs}, Australas. J. Combin. \textbf{59} (2014),
107--119.
\url{https://ajc.maths.uq.edu.au/pdf/59/ajc_v59_p107.pdf}.

\bibitem{JMTVW}
T.~Jiang, D.~Mubayi, Zs.~Tuza, V.~Voloshin and D.~B. West, \emph{The
chromatic spectrum of mixed hypergraphs}, Graphs Combin. \textbf{18}
(2002), 309--318.
\href{https://doi.org/10.1007/s003730200023}{doi:10.1007/s003730200023}.

\bibitem{TV08}
Zs.~Tuza and V.~Voloshin, \emph{Problems and results on colorings of
mixed hypergraphs}, in: Horizons of Combinatorics, Bolyai Soc. Math.
Stud. 17, Springer, 2008, 235--255.
\href{https://doi.org/10.1007/978-3-540-77200-2_12}{doi:10.1007/978-3-540-77200-2\_12}.
Quoted from the authors' copy,
\url{https://spectrum.troy.edu/voloshin/horizon-pf.pdf}.

\bibitem{Voloshin95}
V.~I. Voloshin, \emph{On the upper chromatic number of a hypergraph},
Australas. J. Combin. \textbf{11} (1995), 25--45.

\bibitem{DRATtrim}
N.~Wetzler, M.~J.~H. Heule and W.~A. Hunt, Jr., \emph{DRAT-trim: efficient
checking and trimming using expressive clausal proofs}, in: Theory and
Applications of Satisfiability Testing (SAT 2014), Lecture Notes in
Comput. Sci. 8561, Springer, 2014, 422--429.
\href{https://doi.org/10.1007/978-3-319-09284-3_31}{doi:10.1007/978-3-319-09284-3\_31}.

\end{thebibliography}

\begin{appendices}
\section{The system and its bicolorings}\label{app:data}

The row for $x$ lists the pairs $y~z$ with $x<y<z$ such that
$\{x,y,z\}$ is a block of $S_{39}$; every block appears once, in the row
of its smallest point.

{\small
\noindent
\begin{tabularx}{\textwidth}{@{}r>{\raggedright\arraybackslash}X@{}}
\toprule
$x$ & blocks $\{x,y,z\}$ with $x<y<z$, written $y~z$\\
\midrule
0 & 1~2; 3~6; 4~5; 7~9; 8~15; 10~13; 11~17; 12~21; 14~19; 16~22; 18~20; 23~33; 24~32; 25~26; 27~30; 28~29; 31~37; 34~35; 36~38 \\
1 & 3~5; 4~6; 7~17; 8~9; 10~12; 11~22; 13~16; 14~15; 18~19; 20~21; 23~38; 24~31; 25~32; 26~33; 27~28; 29~37; 30~34; 35~36 \\
2 & 3~4; 5~6; 7~20; 8~14; 9~16; 10~19; 11~21; 12~17; 13~15; 18~22; 23~26; 24~25; 27~35; 28~37; 29~30; 31~32; 33~36; 34~38 \\
3 & 7~21; 8~11; 9~13; 10~15; 12~22; 14~16; 17~18; 19~20; 23~30; 24~37; 25~31; 26~36; 27~32; 28~33; 29~34; 35~38 \\
4 & 7~15; 8~22; 9~20; 10~18; 11~12; 13~17; 14~21; 16~19; 23~36; 24~35; 25~37; 26~29; 27~31; 28~34; 30~32; 33~38 \\
5 & 7~11; 8~17; 9~18; 10~22; 12~16; 13~14; 15~20; 19~21; 23~24; 25~30; 26~28; 27~29; 31~36; 32~35; 33~34; 37~38 \\
6 & 7~22; 8~18; 9~12; 10~14; 11~16; 13~20; 15~19; 17~21; 23~35; 24~30; 25~33; 26~34; 27~38; 28~36; 29~31; 32~37 \\
7 & 8~31; 10~25; 12~26; 13~37; 14~24; 16~33; 18~23; 19~28; 27~36; 29~38; 30~35; 32~34 \\
8 & 10~38; 12~35; 13~33; 16~27; 19~24; 20~34; 21~36; 23~28; 25~29; 26~32; 30~37 \\
9 & 10~26; 11~23; 14~29; 15~35; 17~31; 19~38; 21~37; 22~36; 24~27; 25~34; 28~32; 30~33 \\
10 & 11~34; 16~29; 17~30; 20~32; 21~23; 24~36; 27~37; 28~31; 33~35 \\
11 & 13~35; 14~37; 15~31; 18~38; 19~27; 20~30; 24~26; 25~28; 29~33; 32~36 \\
12 & 13~30; 14~31; 15~36; 18~34; 19~25; 20~27; 23~37; 24~29; 28~38; 32~33 \\
13 & 18~28; 19~34; 21~27; 22~24; 23~32; 25~38; 26~31; 29~36 \\
14 & 17~23; 18~32; 20~38; 22~33; 25~27; 26~30; 28~35; 34~36 \\
15 & 16~32; 17~29; 18~24; 21~33; 22~28; 23~25; 26~37; 27~34; 30~38 \\
16 & 17~37; 18~36; 20~26; 21~25; 23~34; 24~38; 28~30; 31~35 \\
17 & 19~32; 20~28; 22~35; 24~34; 25~36; 26~38; 27~33 \\
18 & 21~29; 25~35; 26~27; 30~31; 33~37 \\
19 & 22~26; 23~29; 30~36; 31~33; 35~37 \\
20 & 22~25; 23~31; 24~33; 29~35; 36~37 \\
21 & 22~30; 24~28; 26~35; 31~34; 32~38 \\
22 & 23~27; 29~32; 31~38; 34~37 \\
\bottomrule
\end{tabularx}
}

\medskip
The $3$-bicoloring has the color classes
\begin{align*}
&\{6,8,11,14,15,17,20,24,27,36\},\\
&\{0,3,10,16,18,19,26,30,32,34,37,38\},\\
&\{1,2,4,5,7,9,12,13,21,22,23,25,28,29,31,33,35\},
\end{align*}
and the $5$-bicoloring is the coloring $g$ of
Section~\ref{sec:construction}, with the classes $\{0\}$, $\{1,2\}$,
$\{3,4,5,6\}$, $\{7,\dots,22\}$ and $\{23,\dots,38\}$.
\end{appendices}

\end{document}
